\documentclass[tikz,border=4pt]{standalone}
\usepackage[utf8]{inputenc}
\usepackage[T1]{fontenc}
\usepackage[brazilian]{babel}
\usepackage{lmodern}
\renewcommand{\familydefault}{\sfdefault}
\usepackage{amsmath}
\usepackage{siunitx}
\sisetup{
  output-decimal-marker = {,},
  range-phrase = { a },
  retain-explicit-plus = true
}
\usetikzlibrary{babel, arrows.meta, positioning, shapes.geometric, fit, calc}

\begin{document}
\begin{tikzpicture}[
  >=Stealth,
  font=\sffamily,
  panelbox/.style={draw=black!40, rounded corners=3pt, fill=gray!3},
  titlebar/.style={draw=black!50, fill=black!8, rounded corners=2pt, font=\small\bfseries, align=center},
  subbox/.style={draw=blue!45!black, fill=blue!7!white},
  pdiff/.style={draw=orange!70!black, fill=orange!24!white},
  oxbox/.style={draw=cyan!60!black, fill=cyan!12!white},
  albox/.style={draw=black!75, fill=gray!38!white},
  polybox/.style={draw=red!55!black, fill=red!22!brown!30!white},
  callout/.style={draw=black!40, fill=white, rounded corners=2pt, align=center},
  bluecallout/.style={draw=blue!65!black, fill=blue!5!white, rounded corners=2pt, align=center},
  pluscharge/.style={circle, draw=red!75!black, fill=red!10!white, inner sep=0.4pt, font=\tiny\bfseries, text=red!80!black}
]

% ==============================================================================
% PAINEL (a): PORTA METÁLICA DE ALUMÍNIO SOBRE SILÍCIO <111>
% ==============================================================================
\draw[panelbox] (0.00, 0.00) rectangle (8.75, 7.25);
\node[titlebar, minimum width=8.75cm, minimum height=0.58cm] at (4.375, 6.96)
  {(a) Porta metálica de alumínio sobre silício $\langle 111 \rangle$};

% Caixa resumo superior (a)
\node[callout, text width=5.18cm, minimum height=1.55cm, font=\fontsize{5.5}{6.8}\selectfont] at (2.85, 5.62) {%
  Alinhamento manual exige folga $L_{ov}$ de \SIrange{3}{5}{\micro\meter}\\
  Fusão do Al (\SI{660}{\celsius}) impede etapas $> \SI{500}{\celsius}$\\
  Sem recozimento em alta T (1 única camada)\\
  $V_{th}$ de \SIrange{-3,5}{-5,0}{\volt} (incompatível com TTL)
};

% Substrato n <111>
\draw[subbox] (0.25, 0.85) rectangle (5.45, 2.50);
\node[font=\scriptsize, text=blue!40!black] at (2.85, 1.10) {Substrato Si tipo $n$ $\langle 111 \rangle$};

% Fonte e Dreno p+ (canal amplo de 2,80 cm entre x=1,45 e x=4,25)
\draw[pdiff] (0.35, 1.68) rectangle (1.45, 2.50);
\node[font=\scriptsize, align=center, text=orange!50!black] at (0.90, 2.09) {Fonte\\[-2pt]($p^+$)};

\draw[pdiff] (4.25, 1.68) rectangle (5.35, 2.50);
\node[font=\scriptsize, align=center, text=orange!50!black] at (4.80, 2.09) {Dreno\\[-2pt]($p^+$)};

% Óxido de campo e óxido fino de porta
\draw[oxbox] (0.25, 2.50) rectangle (0.52, 2.92);
\draw[oxbox] (0.80, 2.50) rectangle (1.02, 2.92);
\draw[oxbox] (4.68, 2.50) rectangle (4.90, 2.92);
\draw[oxbox] (5.18, 2.50) rectangle (5.45, 2.92);
\draw[oxbox] (1.02, 2.50) rectangle (4.68, 2.80);

% Contatos metálicos S e D
\draw[albox] (0.52, 2.50) rectangle (0.80, 3.45);
\draw[thick] (0.66, 3.45) -- (0.66, 3.65);
\node[font=\scriptsize\bfseries] at (0.66, 3.83) {S};

\draw[albox] (4.90, 2.50) rectangle (5.18, 3.45);
\draw[thick] (5.04, 3.45) -- (5.04, 3.65);
\node[font=\scriptsize\bfseries] at (5.04, 3.83) {D};

% Porta de Alumínio com sobreposição mecânica L_ov
\draw[albox] (1.02, 2.80) rectangle (4.68, 3.34);
\node[font=\footnotesize] at (2.75, 3.07) {Porta de Alumínio};

% Destaque da capacitância C_gd na sobreposição do dreno
\draw[draw=blue!70!black, thick, dashed, fill=blue!15!white, fill opacity=0.45]
  (4.25, 2.48) rectangle (4.68, 2.82);

% Cotas L_ov (esquerda e direita)
\draw[dashed, black!60] (1.02, 3.34) -- (1.02, 3.94);
\draw[dashed, black!60] (1.45, 2.50) -- (1.45, 3.94);
\draw[<->, black!80] (1.02, 3.78) -- (1.45, 3.78);
\node[font=\scriptsize] at (1.235, 4.08) {$L_{ov}$};

\draw[dashed, black!60] (4.25, 2.50) -- (4.25, 3.94);
\draw[dashed, black!60] (4.68, 3.34) -- (4.68, 3.94);
\draw[<->, black!80] (4.25, 3.78) -- (4.68, 3.78);
\node[font=\scriptsize] at (4.465, 4.08) {$L_{ov}$};

% Chamada de C_gd parasita elevada (caminho ortogonal sem cruzar textos ou cotas)
\node[bluecallout, font=\scriptsize, inner sep=2.5pt] (cgdnote) at (2.85, 4.26)
  {$C_{gd}$ elevada};
\draw[->, blue!70!black, thick, rounded corners=2pt]
  (cgdnote.east) -- (4.06, 4.26) -- (4.06, 2.64) -- (4.23, 2.64);

% Alta densidade de cargas fixas (+) na interface SiO2/Si<111>
\foreach \x in {1.76, 2.11, 2.46, 2.81, 3.16, 3.51, 3.86} {
  \node[pluscharge] at (\x, 2.50) {+};
}
\node[font=\scriptsize, align=center, text=red!75!black] (acharges) at (2.85, 1.92)
  {Alta carga\\[-2pt]fixa ($+$)};
\draw[->, red!70!black, thin] (2.42, 2.18) -- (2.26, 2.40);
\draw[->, red!70!black, thin] (3.28, 2.18) -- (3.44, 2.40);

% Nota inferior (a)
\node[callout, text width=5.18cm, minimum height=0.58cm, font=\fontsize{6.3}{7.4}\selectfont] at (2.85, 0.42)
  {Cargas ($+$) em $\langle 111 \rangle$ repelem\\[-1pt]lacunas e elevam $|V_{th}|$};

% Mini-inversor ilustrado (Efeito Miller)
\draw[draw=blue!45!black, dashed, rounded corners=3pt, fill=white] (5.60, 0.16) rectangle (8.58, 6.46);
\node[font=\footnotesize\bfseries, align=center, text=blue!55!black] at (7.09, 6.12)
  {Inversor\\[-1pt](Efeito Miller)};

% Trilho V_DD e Carga
\node[font=\scriptsize] at (7.26, 5.56) {$V_{DD}$};
\draw[thick] (7.26, 5.40) -- (7.26, 5.10);
\draw[thick, fill=gray!12] (7.04, 4.42) rectangle (7.48, 5.10);
\node[font=\scriptsize] at (7.26, 4.76) {Carga};
\draw[thick] (7.26, 4.42) -- (7.26, 3.45);

% Saída V_out
\fill (7.26, 4.00) circle (1.4pt);
\draw[thick] (7.26, 4.00) -- (7.74, 4.00);
\node[font=\scriptsize, text=blue!70!black] at (8.10, 4.00) {$V_{out}\!\downarrow$};

% Transistor PMOS chave
\draw[thick] (7.26, 3.45) -- (7.03, 3.45) -- (7.03, 2.55) -- (7.26, 2.55);
\draw[thick] (6.93, 2.65) -- (6.93, 3.35);
\draw[thick] (7.26, 2.55) -- (7.26, 2.08);
\node[font=\scriptsize] at (7.26, 1.90) {$V_{SS}$};

% Entrada V_in
\draw[thick] (6.10, 3.00) -- (6.93, 3.00);
\node[font=\scriptsize, text=red!70!black] at (6.06, 2.72) {$V_{in}\!\uparrow$};

% Capacitor Miller C_gd entre porta e dreno/saída
\fill (6.50, 3.00) circle (1.3pt);
\draw[thick, blue!70!black] (6.50, 3.00) -- (6.50, 3.44);
\draw[very thick, blue!70!black] (6.33, 3.44) -- (6.67, 3.44);
\draw[very thick, blue!70!black] (6.33, 3.58) -- (6.67, 3.58);
\draw[thick, blue!70!black] (6.50, 3.58) -- (6.50, 4.00) -- (7.26, 4.00);
\node[font=\scriptsize\bfseries, text=blue!70!black] at (6.04, 3.51) {$C_{gd}$};

% Explicação concisa Miller
\node[font=\fontsize{6.3}{7.4}\selectfont, align=center, text width=2.76cm] at (7.09, 0.88) {%
  $C_{in,M} = C_{gd}(1 + |A_v|)$\\[2pt]
  Saída em sentido oposto\\[-1pt]amplia carga na entrada
};

% ==============================================================================
% PAINEL (b): PORTA AUTOALINHADA DE POLISSILÍCIO (SGT) SOBRE SILÍCIO <100>
% ==============================================================================
\draw[panelbox] (9.15, 0.00) rectangle (18.30, 7.25);
\node[titlebar, minimum width=9.15cm, minimum height=0.58cm] at (13.725, 6.96)
  {(b) Porta autoalinhada de polissilício, SGT (Si $\langle 100 \rangle$)};

% Banner resumo superior (b): foco em sequência de processo e orçamento térmico
\node[callout, text width=5.20cm, minimum height=1.25cm, font=\fontsize{5.5}{6.8}\selectfont] at (12.00, 5.76) {%
  Polissilício antes da difusão ($> \SI{1100}{\celsius}$)\\
  Permite recozimento em alta T e captura iônica\\
  Óxido contínuo isola 2\textsuperscript{a} camada de Al
};

% Substrato n <100>
\draw[subbox] (9.30, 0.85) rectangle (14.70, 2.50);
\node[font=\scriptsize, text=blue!40!black] at (12.00, 1.10) {Substrato Si tipo $n$ $\langle 100 \rangle$};

% Fonte e Dreno p+ autoalinhados (pequena difusão lateral sob a porta; canal de 2,44 cm)
\draw[pdiff] (9.40, 1.68) rectangle (10.78, 2.50);
\node[font=\scriptsize, align=center, text=orange!50!black] at (10.09, 2.09) {Fonte\\[-2pt]($p^+$)};

\draw[pdiff] (13.22, 1.68) rectangle (14.60, 2.50);
\node[font=\scriptsize, align=center, text=orange!50!black] at (13.91, 2.09) {Dreno\\[-2pt]($p^+$)};

% Óxido isolante intermediário contínuo envolvendo o polissilício e descendo ao substrato lateral
\filldraw[oxbox]
  (9.30, 2.50) -- (9.30, 2.72) -- (10.36, 2.72) -- (10.36, 3.88) --
  (13.64, 3.88) -- (13.64, 2.72) -- (14.70, 2.72) -- (14.70, 2.50) --
  (13.30, 2.50) -- (13.30, 3.36) -- (10.70, 3.36) -- (10.70, 2.50) -- cycle;
\node[font=\scriptsize, text=cyan!50!black] at (12.00, 3.62) {Óxido isolante};

% Óxido fino de porta
\draw[oxbox] (10.70, 2.50) rectangle (13.30, 2.80);

% Porta de polissilício p+
\draw[polybox] (10.70, 2.80) rectangle (13.30, 3.36);
\node[font=\footnotesize, text=red!40!black] at (12.00, 3.08) {Polissilício ($p^+$)};

% Segunda camada condutora independente apoiada integralmente sobre o óxido contínuo
\draw[albox] (10.48, 3.88) rectangle (13.52, 4.42);
\node[font=\footnotesize] at (12.00, 4.15) {2\textsuperscript{a} camada: Alumínio};

% Setas de difusão térmica autoalinhada (todas encerram na mesma cota y=2,76)
\node[font=\scriptsize, text=orange!60!black] at (9.79, 3.74) {Difusão};
\draw[->, orange!75!black, thick, dashed] (9.64, 3.55) -- (9.64, 2.76);
\draw[->, orange!75!black, thick, dashed] (9.94, 3.55) -- (9.94, 2.76);

\node[font=\scriptsize, text=orange!60!black] at (14.21, 3.74) {Difusão};
\draw[->, orange!75!black, thick, dashed] (14.06, 3.55) -- (14.06, 2.76);
\draw[->, orange!75!black, thick, dashed] (14.36, 3.55) -- (14.36, 2.76);

% Destaque C_gd mínima (apenas difusão lateral)
\draw[draw=blue!70!black, thick, fill=blue!20!white, fill opacity=0.5]
  (13.22, 2.48) rectangle (13.32, 2.82);

% Baixa densidade de cargas fixas (+) em <100>
\node[pluscharge] at (11.55, 2.50) {+};
\node[pluscharge] at (12.45, 2.50) {+};
\node[font=\scriptsize, align=center, text=red!75!black] at (12.00, 1.92)
  {Baixa carga\\[-2pt]fixa ($+$)};

% Nota inferior (b)
\node[callout, text width=5.20cm, minimum height=0.58cm, font=\fontsize{6.3}{7.4}\selectfont] at (12.00, 0.42)
  {Interface $\langle 100 \rangle$ reduz cargas ($+$)\\[-1pt]e abaixa $|V_{th}|$};

% Coluna de cartões explicativos à direita em (b): ganhos estruturais e elétricos
\node[callout, text width=2.88cm, minimum height=1.18cm, font=\fontsize{6.5}{7.6}\selectfont] (bcard1) at (16.55, 5.82) {%
  {\bfseries Autoalinhamento}\\[2pt]
  $\Delta L_{align}=0$ limita $L_{ov}$\\[-1pt]à difusão lateral
};

\node[callout, text width=2.88cm, minimum height=1.18cm, font=\fontsize{6.5}{7.6}\selectfont] (bcard2) at (16.55, 4.26) {%
  {\bfseries 2 níveis de fiação}\\[2pt]
  Eliminam túneis difusos\\[-1pt]e reduzem área
};
\draw[->, black!70, thick] (bcard2.west) -- (13.55, 4.26);

\node[bluecallout, text width=2.88cm, minimum height=1.18cm, font=\fontsize{6.5}{7.6}\selectfont] (bcard3) at (16.55, 2.70) {%
  {\bfseries $C_{gd}$ menor (3 a 5$\times$)}\\[2pt]
  Atenua efeito Miller e\\[-1pt]atraso em até 5$\times$
};
\draw[->, blue!70!black, thick] ([yshift=-0.05cm]bcard3.west) -- (13.34, 2.65);

\node[callout, text width=2.88cm, minimum height=1.18cm, font=\fontsize{6.5}{7.6}\selectfont] (bcard4) at (16.55, 1.14) {%
  {\bfseries Substrato $\langle 100 \rangle$ e $p^+$}\\[2pt]
  $V_{th}$ de \SIrange{-1,5}{-2,5}{\volt}\\[-1pt](compatível com TTL)
};
\draw[->, black!70, thick] (bcard4.west) -- (14.72, 1.14);

\end{tikzpicture}
\end{document}
